\section{Admin Web}
\label{sec:AdminWeb}

IPTO includes a web-based administration client under \texttt{apps/admin-web}. This application is not part of the conceptual core of IPTO in the same way as the database schema or the SDL configuration layer, but it gives administrators and developers a browser-based view of the repository.

\subsection{Role in the system}

The web client is a SvelteKit application for:
\begin{itemize}
\item schema administration,
\item tenant-scoped tree browsing,
\item unit search and result inspection.
\end{itemize}

This makes it a client of IPTO rather than a defining part of the persistence model. It sits on top of the repository and API layers described elsewhere in the manual.

\subsection{Structure}

The application is organized around a clear split between routes and reusable UI components:
\begin{description}
\item[\texttt{src/routes/+page.svelte}] Overview page.
\item[\texttt{src/routes/admin/+page.svelte}] Administration views for attributes, records and templates.
\item[\texttt{src/routes/trees/+page.svelte}] Tenant-scoped unit tree browsing.
\item[\texttt{src/routes/search/+page.svelte}] Search UI, saved searches, filters, results and unit detail.
\item[\texttt{src/lib/components}] Shared UI components such as the header, lists, filters and detail panels.
\item[\texttt{src/lib/api.js}] The client-side API access layer.
\end{description}

\subsection{Integration points}

The application is structured to call backend HTTP endpoints rather than embedding repository logic locally. API access goes through \texttt{fetch(...)} calls in \texttt{src/lib/api.js}, which keeps the browser layer focused on interaction and presentation.

The integration surface includes operations related to:
\begin{itemize}
\item tenants,
\item attributes,
\item records,
\item templates,
\item saved searches,
\item unit search and unit detail retrieval.
\end{itemize}

That is consistent with the role of the application: it is an operational and administrative frontend over repository-backed services.

\subsection{Packaging}

The web app can be built into the Quarkus static resources folder so that it is served by the same HTTP server as the Java backend. This is a useful deployment pattern because it reduces the number of independently deployed artifacts for the common case of a Java-based runtime.

\subsection{Observations}

The admin web is a practical systems interface: it exposes repository capabilities to a browser user, but it does so by relying on the backend models and services described in the rest of this manual.
